\chapter{Forskinga utan objekt}

\begin{motto}
Eit fag som har gjort si eiga manglande etterprøvbarheit\\ til ein metode, har slutta å vere eit fag som kan ta feil.
\end{motto}

Då designfaget skulle inn i universitetet for fullt, møtte det eit krav det ikkje kunne oppfylle: krav om forsking. Universitetet er bygd kring frambringing av etterprøvbar kunnskap. For å høyre heime der, måtte designfaget vise at det produserer noko slikt. Det hadde to ærlege val. Det kunne bygd eit fundament som gjorde forskinga mogleg — utvikla presise omgrep, falsifiserbare påstandar, kumulative resultat. Eller det kunne innrømt at det er eit handverksfag og bedt om ein annan plass enn forskingsuniversitetet sin. Det gjorde verken det eine eller det andre. Det fann ein tredje veg: det omdefinerte forsking slik at det det allereie gjorde, kunne telje som forsking. Resultatet er det som går under namnet \emph{forsking gjennom design}\kref{7}, og det er den mest sofistikerte forsvarsmekanismen faget har bygd.

\section{Frayling og dei tre prepossisjonane}

Christopher Frayling sette opp den kanoniske inndelinga\kref{7}: forsking \emph{om} design, forsking \emph{for} design, og forsking \emph{gjennom} design. Dei to fyrste er uproblematiske. Forsking om design er designhistorie, designteori, studiar av faget; det er ekte forsking, men det er som regel historie eller samfunnsvitskap brukt på design, ikkje ein eigen designvitskap. Forsking for design er kunnskap framskaffa for å brukast i designarbeid; det er òg ekte, men det er som regel psykologi, materialvitskap eller marknadsforsking. Den tredje kategorien er den faget hengde si akademiske framtid på, fordi det er den einaste som gjer sjølve designpraksisen til forsking: kunnskap produsert \emph{gjennom} handlinga å designe.

Påstanden er at når ein dugande utøvar designar, og reflekterer over det medan han gjer det, så genererer han ny kunnskap; og at artefaktet han lagar, saman med refleksjonen kring det, utgjer eit forskingsbidrag. Dette er Schön\kref{2} løfta frå pedagogikk til epistemologi. Og her må vi vere presise om kva som måtte vere sant for at det skulle halde. For at å lage eit artefakt skal vere forsking, må artefaktet eller prosessen produsere ein påstand om verda som andre kan ta opp, etterprøve, byggje vidare på, eller vise er feil. Det er minstekravet. Ikkje at det skal vere som naturvitskap; berre at det skal vere kumulativt og etterprøvbart, slik at feltet kan lære noko som varer ut over den einskilde forskaren.

\section{Kvifor det ikkje held}

Sjå då på kva eit typisk forsking-gjennom-design-arbeid faktisk leverer. Ein forskar designar eit artefakt — eit produkt, ein teneste, ein installasjon, ein spekulativ gjenstand. Forskaren reflekterer over prosessen og skriv ein tekst om kva han lærte. Artefaktet og teksten saman vert lagde fram som bidraget. Spør no: kva er påstanden? Kva ville telje som å motbevise han? Kan ein annan forskar etterprøve han? I det overveldande fleirtalet av tilfelle er svaret at det ikkje finst nokon avgrensa påstand, at ingenting ville telje som motbevis, og at ingen kan etterprøve noko fordi resultatet er bunde til den einskilde forskaren sin praksis, kontekst og person. Arbeidet er ofte rikt, gjennomtenkt, til og med vakkert. Men det er ikkje kumulativt. Den neste forskaren byrjar ikkje der den førre slutta; ho byrjar på nytt, med sitt eige artefakt, sin eigen refleksjon. Etter tretti år med forsking gjennom design finst det ingen veksande kropp av etablerte resultat som studentane lærer som grunnlag, slik biologistudenten lærer seleksjon eller fysikkstudenten lærer mekanikk. Det finst ein veksande bunke av einskildstudiar som ikkje snakkar saman.

Og no kjem grepet som gjer det heile uangripeleg. I staden for å sjå denne mangelen på kumulativitet som ein veikskap, har faget gjort han til ein dyd, ved hjelp av eit lånt omgrep: \emph{vrange problem}. Rittel og Webber innførte omgrepet for planlegging\kref{4}; vrange problem er problem utan endeleg formulering, utan stoppregel, utan rett-eller-gale-svar, der kvar løysing er eit eingongstilfelle. Buchanan importerte det til design\kref{5} og gjorde det til ein hjørnestein. Og lat oss vere rettferdige: mange designproblem \emph{er} vrange i denne forstanden. Innvendinga er ikkje at omgrepet er gale. Innvendinga er korleis faget brukar det. Det brukar «vrange problem» til å forklare kvifor forskinga ikkje treng vere kumulativ, kvifor resultata ikkje treng vere etterprøvbare, kvifor det ikkje finst rette svar å konvergere mot. Med andre ord: faget tek det som skulle vore objektet for ein teori — det faktum at formgjeving navigerer mange motstridande krav samstundes — og brukar det som ei orsaking for å sleppe å lage teorien.

Dette er den djupaste feilen, så lat meg seie han skarpt. At eit problem er vrangt, tyder ikkje at det ikkje kan finnast ein teori om kvifor det er vrangt. Tvert om: at formgjeving må balansere mange samtidige, motstridande seleksjonstrykk over eit høgdimensjonalt formrom, er ikkje ein grunn til at teori er umogleg; det er \emph{innhaldet} i den teorien som trengst. Biologien sine problem er vrange i nett same forstand — ein organisme må løyse utalige motstridande krav samstundes — og biologien svara ikkje med å erklære seg sjølv uvitskapleg. Han bygde tilpassingslandskapet, eit presist matematisk apparat for nett denne situasjonen\kref{12}. Designfaget hadde same val og tok det motsette. Det tok vrangskapen som bevis på at det ikkje kunne ha eit fundament, i staden for som spesifikasjonen på kva fundamentet måtte handtere. «Vrange problem» vart faget si vitskapsteoretiske immunforsvarscelle: kvar gong nokon spør etter etterprøvbar kunnskap, vert han mint på at problema er vrange, og difor er kravet hans malplassert. Eit fag som har immunisert seg mot kravet om å kunne ta feil, har ikkje løyst problemet med å vere ein vitskap. Det har berre slutta å vere eit fag som kan ta feil — og det som ikkje kan ta feil, kan heller ikkje lære.

\section{Doktorgraden som ritual}

Konsekvensen ser du i den praksisbaserte doktorgraden. Ein kandidat brukar år på eit kreativt arbeid, supplerer det med ein refleksiv tekst, og får tittelen doktor. Eg seier ikkje at arbeida er dårlege; mange er framifrå som arbeid. Eg seier at tittelen er lausriven frå det han er meint å sertifisere. Ein doktorgrad sertifiserer at innehavaren har produsert eit etterprøvbart, originalt bidrag til ein kumulativ kunnskapskropp. Når kunnskapskroppen ikkje er kumulativ og bidraget ikkje er etterprøvbart, sertifiserer tittelen ingenting av dette. Han sertifiserer at kandidaten gjorde eit seriøst kreativt arbeid og reflekterte godt over det. Det er ein heider verdt å gje — men det er ein annan heider enn den tittelen lovar, og faget brukar med vilje den gamle tittelen for den nye tingen, fordi den gamle tittelen ber den autoriteten faget vil ha. Det er den same dobbeltbokføringa som i kapittel \textsc{i}: faget krev forskaren sin status for handverkaren sitt arbeid.

\vignett

Faget kan altså ikkje grunngje dommane sine, ikkje fordøye låna sine, og ikkje gjere forskinga si kumulativ. Ein kunne framleis halde at dette er akademiske spørsmål utan praktiske konsekvensar — at faget fungerer i praksis trass i at det ikkje kan grunngje seg sjølv. Det neste kapitlet viser at det ikkje stemmer. Mangelen på fundament har ein pris, og prisen vert betalt av andre enn designarane.
